\documentclass[a4paper,12pt]{article}
\usepackage{graphicx}       % For including images
\usepackage{amsmath}        % For mathematical symbols
\usepackage{listings}       % For code listings
\usepackage{xcolor}         % For color in code listings
\usepackage{hyperref}       % For clickable links
\usepackage{caption}        % For captions
\usepackage{subcaption}     % For subfigures
\usepackage{fancyhdr}       % For headers and footers
\usepackage{geometry}       % For page margins
\usepackage{titling}        % For title customizations
\usepackage{tocbibind}      % For including TOC
\usepackage[dvipsnames]{xcolor}
\usepackage{float} 
\usepackage{longtable}      % For long tables
\usepackage{booktabs}       % For professional tables
\usepackage{enumitem}       % Beeter lists

% --- Code Listing Configuration ---
\lstset{
    language=C++,
    basicstyle=\ttfamily\small,
    keywordstyle=\color{blue}\bfseries,
    commentstyle=\color{OliveGreen},
    stringstyle=\color{Red},
    numbers=left,
    numberstyle=\tiny\color{gray},
    stepnumber=1,
    showstringspaces=false,
    tabsize=4,
    breaklines=true,
    breakatwhitespace=true,
    frame=single,
    captionpos=b
}

% Custom Keywords for Project Classes
\lstset{
    classoffset=1, 
    morekeywords={ProcessInfo, ProcessScanner, TUI, RuleEngine, ProcessController, CommandParser, CommandType, Command, Warning, vector, map, string, ifstream}, 
    keywordstyle=\color{TealBlue}\bfseries, 
    classoffset=0 
}

% Margins
\geometry{top=2.5cm, bottom=2.5cm, left=2.5cm, right=2.5cm}

\begin{document}

% --- Check Title Page ---
\thispagestyle{empty}
\vspace*{1cm}
\begin{center}
    {\Huge\textbf{Advanced Linux Process Manager}} \\
    \vspace{0.5cm}
    {\Huge\textbf{System Programming Lab Report}}
    \vspace{1.5cm}

    % \includegraphics[width=0.3\textwidth]{logo.png} % User can add logo here
    
    {\large Session 2023-2027}\\[0.2cm]
    \vspace{1cm}
    {\Large\textbf{Submitted By:}}\\[0.6cm]
    {\large Mohammad Bilal (2023-CS-168)}\\[1.5cm]
    
    {\Large\textbf{Supervised By:}}\\[0.6cm]
    {\large Mr. Nazeef Ul Haq}\\
    {\large Mr. Waseem}\\[2cm]
    
    {\LARGE \textbf{Department of Computer Science}}\\[0.5cm]
    {\Huge \textbf{University of Engineering and Technology, Lahore}}\\[0.5cm]
    \vfill
\end{center}

\newpage
\pagestyle{fancy}
\fancyhf{}
\fancyhead[L]{\textbf{Process Manager Project}}
\fancyhead[R]{OS Lab}
\fancyfoot[C]{\thepage}
\fancyfoot[R]{2023-CS-168}

\newpage
\tableofcontents
\newpage
% \listoffigures
% \newpage

% --- Section 1: Introduction ---
\section{Introduction}

\subsection{Project Overview}
The \textbf{Linux Process Manager} is a sophisticated, terminal-based system utility designed to monitor and control processes in a Linux environment (including WSL). Developed in C++17, effectively mimicking and extending functionality found in standard tools like \texttt{top} or \texttt{htop}.

Unlike standard utilities, this project includes an \textbf{AI Rule Engine} that actively monitors process behavior over time. Instead of simply displaying static numbers, the system maintains a historical buffer (60-second rolling window) for every process. This allows it to detect complex anomalies such as memory leaks, zombie processes, and CPU spikes, providing actionable intelligence to the user.

\subsection{Objective & Scope}
The core objectives of this project were:
\begin{itemize}
    \item \textbf{System Call Proficiency}: Mastery of Linux systems programming interfaces (\texttt{ptrace}, \texttt{kill}, \texttt{sysconf}).
    \item \textbf{File System Parsing}: Deep understanding of the Linux \texttt{/proc} pseudo-filesystem.
    \item \textbf{Data Structures}: Implementation of custom structures for tracking history and efficient map-based lookups.
    \item \textbf{UI Design}: Creating a responsive, flicker-free Text User Interface (TUI) using \texttt{ncurses}.
    \item \textbf{Algorithmic Logic}: Developing heuristics for the "AI" component to analyze resource trends.
\end{itemize}

\subsection{Feature Summary}
The application supports a rich set of features:
\begin{enumerate}
    \item \textbf{Real-Time Monitoring}: Displays PID, Name, State, Thread counts, CPU\%, Memory (RSS), and Disk I/O speeds.
    \item \textbf{Historical Analysis}: Tracks resource usage over the last minute to generate trend graphs.
    \item \textbf{Process Control}: Supports \texttt{kill}, \texttt{pause} (SIGSTOP), \texttt{resume} (SIGCONT), and \texttt{priority} adjustment.
    \item \textbf{Visualization}: Includes a hierarchical \texttt{tree} view and detailed single-process inspection screens.
    \item \textbf{AI Safety Guard}: Warns users about potentially malicious or buggy processes (Zombies, Leaks).
\end{enumerate}

\newpage

% --- Section 2: Comparison ---
\section{Comparison with Existing Tools}

The custom Process Manager was designed to address specific educational and usability gaps found in traditional tools.

\subsection{Comparison Matrix}

\begin{table}[H]
    \centering
    \caption{Feature Comparison with Industry Standards}
    \begin{tabular}{|l|p{4cm}|p{4cm}|p{4cm}|}
    \hline
    \textbf{Feature} & \textbf{Standard 'top'} & \textbf{'htop'} & \textbf{Our Project} \\ \hline
    \textbf{Interactivity} & Limited (shortcuts) & Good (Mouse/Key) & \textbf{Command-Line Prompt} \\ \hline
    \textbf{History} & None (Instant only) & None & \textbf{60s Rolling Buffer} \\ \hline
    \textbf{Analysis} & Passive (User must watch) & Passive & \textbf{Active AI Alerts} \\ \hline
    \textbf{Visualization} & Raw Text & Bars/Colors & \textbf{ASCII Graphs \& Tree} \\ \hline
    \textbf{Structure} & Monolithic C & Structured C & \textbf{Modular OOP C++} \\ \hline
    \end{tabular}
\end{table}

\subsection{Key Differentiators}
\begin{itemize}
    \item \textbf{Command Prompt Paradigm}: While \texttt{top} relies on single-key shortcuts (e.g., press 'k' then type PID), our tool offers a persistent command line at the bottom. This allows for more expressive commands like `filter chrome`, `sort cpu`, or `ai 1234 ignore`. This is more intuitive for users comfortable with shells.
    \item \textbf{Automated Diagnostics}: A novice user looking at `top` might not notice a process slowly consuming 1MB more RAM every second. Our **Rule Engine** explicitly calculates the regression slope of memory usage and flashes a `MEMORY_LEAK` warning in red, teaching the user what to look for.
    \item \textbf{Zombie Reaping}: Our tool doesn't just display zombies; it includes logic to actively identify them and, if they are children of the manager itself, clean them up automatically.
\end{itemize}

\newpage

% --- Section 3: Architecture ---
\section{System Architecture}

The project follows a rigorous Object-Oriented design, separating concerns into Acquisition, Storage, Analysis, and Presentation layers.

\subsection{Class Diagram Description}
\begin{itemize}
    \item \textbf{ProcessInfo}: The fundamental data unit. It acts as a Data Transfer Object (DTO) containing all stats for a single PID. It owns the `std::vector` history buffers.
    \item \textbf{ProcessScanner}: The "Driver" layer. It interfaces with the OS Kernel via `/proc`. It is responsible for raw file parsing and delta calculations (converting ticks to percentages).
    \item \textbf{ProcessController}: The "Executor" layer. It abstracts standard C library system calls. It ensures safety (e.g., checking for valid PIDs) before executing destructive actions.
    \item \textbf{RuleEngine}: The "Brain" layer. It is a pure logic component that takes a read-only map of processes and returns a vector of `Warning` objects. It is stateless between runs, relying on the process history.
    \item \textbf{TUI}: The "View" layer. It knows nothing about system calls. It simply renders the data it is given using `ncurses` primitives.
\end{itemize}

\subsection{Data Flow}
1.  **Tick (100ms)**: The main loop wakes up.
2.  **Scan**: `ProcessScanner` reads hundreds of files in `/proc`.
3.  **Update**: It matches existing PIDs to the previous map. If a match is found, it calculates `(Current - Prev)` to derive CPU\% and I/O speeds.
4.  **Analyze**: The map is passed to `RuleEngine`. Rules are applied (e.g., `check_memory_leak`).
5.  **Render**: The map + warnings are passed to `TUI`. The table is drawn.
6.  **Input**: User keystrokes are captured. if `Enter` is pressed, `CommandParser` routes the action to `ProcessController`.

\newpage

% --- Section 4: Deep Dive Features ---
\section{Detailed Feature Analysis}

\subsection{1. Advanced Process Discovery}
The `ProcessScanner` does not simply list files. It performs complex correlation.
\begin{itemize}
    \item \textbf{Command Line Reconstruction}: It reads `/proc/[pid]/cmdline`, which is null-delimited. The separate arguments are joined to show the full launch command.
    \item \textbf{User/Parent Mapping}: It reads the `PPID` to build the hierarchy used in the `Tree` view.
\end{itemize}

\subsection{2. Intelligent Resource Calculation}
Calculating accurate CPU usage in Linux is non-trivial.
\begin{enumerate}
    \item Read process usage: \texttt{utime} (user mode ticks) + \texttt{stime} (kernel mode ticks) from \texttt{/proc/[pid]/stat}.
    \item Read system usage: Sum of all CPU states (user, nice, system, idle, iowait) from \texttt{/proc/stat}.
    \item **Formula**:
    $$ \text{Process} \Delta = (\text{Total}_{\text{now}} - \text{Total}_{\text{prev}}) $$
    $$ \text{System} \Delta = (\text{SystemTotal}_{\text{now}} - \text{SystemTotal}_{\text{prev}}) $$
    $$ \text{Usage \%} = \left( \frac{\text{Process} \Delta}{\text{System} \Delta} \right) \times \text{Core Count} \times 100 $$
\end{enumerate}
This handling of multi-core normalization ensures that a single thread running at 100\% shows as "100\%" (or "25\%" on a 4-core system depending on `top` mode preferences, we normalized to core count).

\subsection{3. Rule-Based AI Engine}
The AI is the standout feature. It implements four specific heuristics:

\subsubsection{Rule A: The Zombie Hunter}
Zombies are dead processes waiting for their parent.
\begin{lstlisting}[language=C++]
void RuleEngine::check_stuck_process(const ProcessInfo& p, vector<Warning>& warnings) {
    if (p.state == "Z") {
        warnings.push_back({p.pid, "ZOMBIE", "Parent waiting..."});
    }
}
\end{lstlisting}

\subsubsection{Rule B: High CPU Detection}
A hard threshold is insufficient because short bursts are normal. The engine checks the \textbf{Average} of the last 30 samples.
\begin{lstlisting}[language=C++]
// Conceptual Logic
double avg = calculate_mean(p.cpu_history, 30); // Last 30 seconds
if (avg > 80.0) {
    warnings.push_back({p.pid, "HIGH_CPU", "Sustained Loop Detected"});
}
\end{lstlisting}

\subsubsection{Rule C: Memory Leak Heuristic}
This detects "Combustive Growth". The algorithm checks if the memory usage has strictly increased or stayed same (non-decreasing) for 60 seconds, with a net positive growth > 10MB.
\begin{lstlisting}[language=C++]
bool is_leaking = true;
for (size_t i = 1; i < p.mem_history.size(); ++i) {
    if (p.mem_history[i] < p.mem_history[i-1]) {
        is_leaking = false; // Usage dropped, probably Garbage Collection
        break;
    }
}
if (is_leaking && (total_growth > 10_MB)) warn(p);
\end{lstlisting}

\newpage

% --- Section 5 Implementation ---
\section{Implementation \& Code Structure}

\subsection{Header Files Overview}
\begin{itemize}
    \item \texttt{ProcessInfo.h}: Defines the structure of a process entry. Changes here propagate to all modules.
    \item \texttt{CommandParser.h}: Defines the grammar of the user command shell.
    \item \texttt{TUI.h}: Encapsulates the ncurses library interactions.
\end{itemize}

\subsection{Key Algorithms}

\subsubsection{Sorting Processes}
Sorting is dynamic based on user selection. We use C++ lambdas for flexibility. This allows extending sorting to any field without rewriting logic.

\begin{lstlisting}[language=C++]
std::sort(procs.begin(), procs.end(), [mode](auto& a, auto& b) {
    if (mode == SortMode::CPU) return a.cpu_percent > b.cpu_percent;
    if (mode == SortMode::MEM) return a.rss > b.rss;
    return a.pid < b.pid;
});
\end{lstlisting}

\subsubsection{Tree View Recursion}
The `tree` command requires transforming the flat process map into a hierarchical structure. We perform this in $\mathcal{O}(N)$ by building an adjacency list first.

\begin{lstlisting}[language=C++]
// 1. Build Adjacency List
map<int, vector<int>> children;
for (auto& p : processes) children[p.ppid].push_back(p.pid);

// 2. Recursive Render
void draw_tree(int pid, int depth) {
    print_indent(depth);
    print_process(pid);
    for (int child : children[pid]) {
        draw_tree(child, depth + 1);
    }
}
\end{lstlisting}

\newpage

% --- Section 6: User Manual ---
\section{User Manual \& Command Reference}

This section serves as a comprehensive guide for operating the shell.

\subsection{Basic Workflow}
1.  **Launch**: Run `./process_manager`.
2.  **Observe**: Watch the main list. Warnings appear in red at the bottom.
3.  **Command**: Type commands in the prompt footer.

\subsection{Visual Navigation Commands}

\subsubsection{Sorting}
Organize the chaos. Essential for finding resource hogs.
\begin{itemize}
    \item \texttt{sort cpu}: Descending order of CPU usage. Best for finding frozen apps.
    \item \texttt{sort mem}: Descending order of RAM. Best for finding leaks.
    \item \texttt{sort pid}: Ascending order of ID. Standard view.
\end{itemize}

\subsubsection{Filtering}
Real-time fuzzy search.
\begin{itemize}
    \item \texttt{filter <query>}: Only shows rows where Name contains \texttt{<query>}.
    \item \textbf{Example}: \texttt{filter python} usually shows all ML scripts running.
    \item \textbf{Reset}: Type \texttt{filter} with no arguments to clear.
\end{itemize}

\subsubsection{Tree View}
\begin{itemize}
    \item \texttt{tree}: Toggles strict parent-child visualization.
    \item \textbf{Usage}: Extremely useful for seeing which processes were spawned by a shell script or server daemon (e.g., `nginx` spawning workers).
\end{itemize}

\subsubsection{Deep Inspection}
\begin{itemize}
    \item \texttt{info <pid>}: Clears the screen and performs a "Deep Dive".
    \item \textbf{Displays}:
        \begin{itemize}
            \item Exact Memory usage in bytes.
            \item System vs User time ratio.
            \item \textbf{ASCII Chart}: A visual bar chart of CPU usage over the last 60 seconds.
        \end{itemize}
\end{itemize}

\subsection{Process Control Commands}
\textit{Warning: These commands affect running software.}

\begin{longtable}{|p{3cm}|p{4cm}|p{7cm}|}
\hline
\textbf{Command} & \textbf{Parameters} & \textbf{Detailed Action} \\ \hline
\texttt{run} & \texttt{<cmd> [args]} & Uses \texttt{fork()} and \texttt{execvp()} to launch a process. \newline Example: \texttt{run sleep 60} creates a background sleeper. \\ \hline
\texttt{kill} & \texttt{<pid>} & Sends \texttt{SIGTERM (15)}. Gives the process a chance to cleanup files. \\ \hline
\texttt{pause} & \texttt{<pid>} & Sends \texttt{SIGSTOP (19)}. Forces the OS to unschedule the thread. Good for pausing CPU hogs temporarily. \\ \hline
\texttt{resume} & \texttt{<pid>} & Sends \texttt{SIGCONT (18)}. Returns a paused process to the run queue. \\ \hline
\texttt{priority} & \texttt{<pid> <val>} & Calls \texttt{setpriority()}. \newline Val: -20 (High Priority) to 19 (Low). \\ \hline
\end{longtable}

\subsection{AI Interaction Commands}
The AI engine is interactive.
\begin{itemize}
    \item \texttt{ai <pid> ignore}: Adds the PID to a \texttt{std::set} of ignored IDs. The RuleEngine skips these, silencing warnings.
    \item \texttt{ai <pid> kill}: A "macro" command that immediately terminates the offender and acknowledges the warning.
\end{itemize}

\newpage

% --- Section 7: Testing ---
\section{Testing & Validation Strategy}

To ensure robustness, specific test scenarios were developed.

\subsection{Test Case 1: The Fork Bomb (Stress Test)}
\textbf{Objective}: Ensure the scanner doesn't crash when thousands of processes appear instantly.
\textbf{Action}: Triggered a rapid process generator.
\textbf{Result}: The TUI slowed down due to map insertions but remained stable. The table expanded correctly.

\subsection{Test Case 2: The Zombie Factory}
\textbf{Objective}: Validate Rule A (Zombie Detection).
\textbf{Script Used}:
\begin{lstlisting}[language=bash]
# A script that creates a child and exits the parent
(sleep 1 & exec /bin/sleep 5)
\end{lstlisting}
\textbf{Observation}: The subprocess becomes orphaned. The Process Manager highlights it with a `[Z]` state and moves it to the bottom or triggers a warning.

\subsection{Test Case 3: Memory Leaker}
\textbf{Objective}: Validate Rule C (Leak Detection).
\textbf{Mechanism}: A C++ program allocating 1MB chunks in a loop without freeing.
\begin{lstlisting}[language=C++]
while(1) { malloc(1024*1024); sleep(1); }
\end{lstlisting}
\textbf{Result}: After approx 45 seconds, the "MEMORY_LEAK" warning appeared in the footer, validating the slope calculation logic.

\subsection{Test Case 4: Permission Handling}
\textbf{Objective}: Interacting with root processes.
\textbf{Action}: Attempted `kill 1` (init).
\textbf{Result}: System call failed with `EPERM`. The Process Controller caught the error and displayed "Permission Denied" in the TUI instead of crashing.

\section{Future Enhancements}

While the system is fully functional, future iterations could include:
\begin{itemize}
    \item \textbf{Network I/O}: Parsing \texttt{/proc/net/dev} to show bandwidth per process.
    \item \textbf{GPU Monitoring}: Integrating \texttt{nvidia-smi} bindings for AI workloads.
    \item \textbf{Daemon Mode}: Running as a background service that logs anomalies to a file instead of displaying them in TUI.
\end{itemize}

\section{Conclusion}

The Linux User-Space Process Manager project provided deep insights into the OS kernel's operation. By eschewing standard libraries in favor of raw \texttt{/proc} parsing, we gained a granular understanding of process states, scheduling, and memory management. The addition of the AI engine transforms the tool from a passive monitor to an active system guardian, demonstrating the potential for intelligent OS utilities.

\end{document}
